% !TeX program = xelatex
\documentclass[12pt,a4paper]{article}
\usepackage[margin=2cm]{geometry}
\usepackage{fontspec}
\usepackage{amsmath,amssymb,enumitem}
\setmainfont{Times New Roman}
\setlength{\parindent}{0pt}
\setlength{\parskip}{6pt}
\begin{document}
\begin{center}
{\Large\bfseries Đề thi tuyển sinh vào lớp 10 Hà Nội năm 2019--2020 -- phần Đại số}
\end{center}

\textbf{Câu 1.} Cho hai biểu thức
\[
A=\frac{4(\sqrt{x}+1)}{25-x},\qquad
B=\left(\frac{15-\sqrt{x}}{x-25}+\frac{2}{\sqrt{x}+5}\right):\frac{\sqrt{x}+1}{\sqrt{x}-5},
\quad x\ge0,\ x\ne25.
\]
\begin{enumerate}[label=\alph*)]
  \item Tính giá trị của biểu thức $A$ khi $x=9$.
  \item Rút gọn biểu thức $B$.
  \item Tìm tất cả các giá trị nguyên của $x$ để biểu thức $P=A\cdot B$ đạt giá trị nguyên lớn nhất.
\end{enumerate}

\textbf{Câu 2.}
\begin{enumerate}[label=\alph*)]
  \item Giải bài toán bằng cách lập phương trình hoặc hệ phương trình.

  Hai đội công nhân cùng làm chung một công việc thì sau $15$ ngày làm xong. Nếu đội thứ nhất làm riêng trong $3$ ngày rồi dừng lại và đội thứ hai làm tiếp công việc đó trong $5$ ngày thì cả hai đội hoàn thành được $25\%$ công việc. Hỏi mỗi đội làm riêng thì bao nhiêu ngày mới hoàn thành xong công việc trên.
  \item Một bồn nước inox có dạng hình trụ với chiều cao $1{,}75$ m và diện tích đáy là $0{,}32$ m$^2$. Hỏi bồn nước này đựng đầy được bao nhiêu mét khối nước? (Bỏ qua bề dày của bồn nước.)
\end{enumerate}

\textbf{Câu 3.}
\begin{enumerate}[label=\arabic*)]
  \item Giải phương trình
  \[
  x^4-7x^2-18=0.
  \]
  \item Trong mặt phẳng tọa độ $Oxy$, cho đường thẳng $(d):y=2mx-m^2+1$ và parabol $(P):y=x^2$.
  \begin{enumerate}[label=\alph*)]
    \item Chứng minh $(d)$ luôn cắt $(P)$ tại hai điểm phân biệt.
    \item Tìm tất cả giá trị của $m$ để $(d)$ cắt $(P)$ tại hai điểm phân biệt có hoành độ $x_1,x_2$ thỏa mãn
    \[
    \frac{1}{x_1}+\frac{1}{x_2}=\frac{-2}{x_1x_2}+1.
    \]
  \end{enumerate}
\end{enumerate}

\textbf{Câu 5.} Cho biểu thức $P=a^4+b^4-ab$ với $a,b$ là các số thực thỏa mãn $a^2+b^2+ab=3$. Tìm giá trị lớn nhất và giá trị nhỏ nhất của $P$.
\end{document}
